\documentclass[12pt,a4paper]{article}
% ------------------------------------------------------
% Idioma y codificación
% ------------------------------------------------------
\usepackage[spanish,es-nodecimaldot]{babel}
\usepackage[utf8]{inputenc}
\usepackage[T1]{fontenc}

% ------------------------------------------------------
% Tipografía y formato
% ------------------------------------------------------
\usepackage{lmodern}    % Fuente Latin Modern
\usepackage{microtype}  % Mejor espaciado
\usepackage{xcolor}     % Colores
\setlength{\parindent}{0pt} %eliminar sangría
\definecolor{UGblue}{RGB}{0,33,71}
\definecolor{UGgold}{RGB}{212,175,55}
\usepackage{graphicx}
\usepackage{float}
\usepackage{setspace}
\usepackage{titlesec}
\usepackage{amsmath}
\usepackage{tikz}
\usetikzlibrary{calc,decorations.pathreplacing,angles,quotes}
\usepackage{pgfplots}
\pgfplotsset{compat=1.18}
\usepackage{caption}
\usepackage{multicol} %Permite hacer columnas de texto

% ------------------------------------------------------
% Matemáticas
% ------------------------------------------------------
\usepackage{amsmath,amssymb,amsthm,mathtools}
\usepackage{mathrsfs}
\usepackage{stmaryrd}
%\usepackage{physics}

% ------------------------------------------------------
% Otros paquetes útiles
% ------------------------------------------------------
\usepackage{enumitem}
\usepackage{csquotes}
\usepackage{hyperref}
\hypersetup{
    colorlinks=true,
    linkcolor=UGblue,
    urlcolor=UGgold,
    citecolor=UGblue
}
\usepackage{geometry}
\geometry{margin=2.5cm}
\usepackage{multicol}

% ------------------------------------------------------
% Teoremas, definiciones, etc.
% ------------------------------------------------------
\theoremstyle{plain}
\newtheorem{teo}{Teorema}[section]
\newtheorem{prop}[teo]{Proposición}
\newtheorem{lema}[teo]{Lema}
\theoremstyle{definition}
\newtheorem{defi}[teo]{Definición}
\theoremstyle{remark}
\newtheorem{ejemplo}[teo]{Ejemplo}
\newtheorem{obs}[teo]{Observación}

% ------------------------------------------------------
% Comandos útiles
% ------------------------------------------------------
\renewcommand{\P}{\mathbb{P}}
\newcommand{\Var}{\text{Var}}
\newcommand{\E}{\mathbb{E}}
\newcommand{\N}{\mathbb{N}}
\newcommand{\Z}{\mathbb{Z}}
\newcommand{\Q}{\mathbb{Q}}
\newcommand{\I}{\mathbb{I}}
\newcommand{\R}{\mathbb{R}}
\newcommand{\C}{\mathbb{C}}
\newcommand{\B}{\mathbb{B}}
\newcommand{\x}{\mathrm x}
\newcommand{\y}{\mathrm y}
\newcommand{\Int}{\mathrm{int}}
\newcommand{\Ext}{\mathrm{ext}}
\newcommand{\diag}{\mathrm{diag}}
\newcommand{\norm}[1]{\left\Vert#1\right\Vert}
\newcommand{\llaves}[1]{\left\{#1\right\}}
\newcommand{\parentesis}[1]{\left(#1\right)}
\newcommand{\corchetes}[1]{\left[#1\right]}
\newcommand{\abs}[1]{\left|#1\right|}
\newcommand{\qedt}{\hfill $\blacktriangle$}
\renewcommand{\qed}{\hfill $\blacksquare$}



\begin{document}
\textsc{\textbf{Tarea 2 Probabilidad.}}\\
    \textsc {Profesor Antonio Murillo Salas.\\
    Hecho por Luis David uicab Martínez.\\}
    09 de septiembre de 2026\\
    \rule{15cm}{4pt}\\[0.3cm]
    
\section*{Ejercicio 1.}
Considere la siguiente función
\[
f(x) =
\begin{cases}
    c(4x - x^2), & \text{si } 0 < x < 4,\\
    0, & \text{en otro caso}.
\end{cases}
\]
\begin{enumerate}
    \item[(i)] Determine el valor de la constante $c$ de modo que $f$ sea una función de densidad.
    \item[(ii)] Sea $X$ la variable aleatoria con función de densidad $f$ (con el valor de $c$ encontrado en (i)). Encuentre la función de distribución de $X$.
    \item[(iii)] Determine la media y la varianza de $X$.
\end{enumerate}
\textbf{Solución.}
\begin{enumerate}
    \item[(i)] Notemos que
        \begin{align*}
            \int_{-\infty}^{\infty}f(x)dx &= \int_{0}^{4}c\left(4x-x^2\right)dx \\
            &= c\left(4\cdot \frac{x^2}{2}-\frac{x^3}{3}\right)\big|_{x=0}^{4} \\
            &= c\cdot\frac{32}{3}.
        \end{align*}
    \item[(ii)] Para la función de distribución, calculamos la siguiente integral
        \begin{align*}
            F_X(x) &= \int_{-\infty}^{x}f(t)dt \\
            &= \int_{0}^{x}\frac{3}{32}\left(4t-t^2\right)dt \\
            &= \frac{3}{32}\int_{0}^{x}\left(4\cdot \frac{t^2}{2}-\frac{t^3}{3}\right)\big|_{t=0}^{x} \\
            &= \frac{6x^2-x^3}{32},
        \end{align*}
        dónde $x\in (0,4)$. Si $x<0$, $F_X(x)=0$ y si $x>4$, $F_X(x)=1$.
    \item[(iii)] Tenemos que
        \begin{align*}
            \E\left(X\right) &= \int_{-\infty}^{\infty}xf(x)dx \\
            &= \frac{3}{32}\int_{0}^{4}\left(4x^2-x^3\right)dx \\
            &= \frac{3}{32}\left(4\cdot\frac{x^3}{3}-\frac{x^4}{4}\right)\Big|_{x=0}^{4} \\
            &= 2.
        \end{align*}
        Por otra parte,
        \begin{align*}
            \E\left(X^2\right) &= \int_{-\infty}^{\infty}x^2f(x)dx \\
            &= \frac{3}{32}\int_{0}^{4}\left(4x^3-x^4\right)dx \\
            &= \frac{3}{32}\left(4\cdot\frac{x^4}{4}-\frac{x^5}{5}\right)\Big|_{x=0}^4 \\
            &= 4.8
        \end{align*}
        En consecuencia,
        \[
        \text{Var}(X) = \E(X^2)-\E(X)^2 = 4.8 - 2^2 = 0.8
        \]
        \hfill $\blacktriangle$
\end{enumerate}
Puesto que queremos que $f$ sea función de distribución, la integral anterior ha de valor 1, y en consecuencia, $c=3/32$.


\section*{Ejercicio 2.}
Sea $X$ una variable aleatoria continua con función de distribución $F$ y sea $p \in (0,1)$. El \textit{percentil} $100$ \textit{p-ésimo} (o percentil $100p$), denotado por $P_p$, se define como el valor que satisface
\[
\P(X \leq P_p) = p.
\]
Si $X$ es una variable aleatoria exponencial con parámetro (tasa) $\lambda > 0$, determine una expresión para $P_p$ con $p \in (0, 1)$.\\
\textbf{Solución.} \\
Por hipótesis, sabemos que $X$ tiene función de densidad
\[
f(x)=
\begin{cases}
    \int_{-\infty}^{x}\lambda e^{-\lambda x}dx, &\text{ si } x > 0 \\
    0, &x\leq 0
\end{cases}
\]
Así, tomando $u=-\lambda t$, se tiene que $dt=-du/\lambda$ y en consecuencia
\begin{align*}
    F_X(x) &= \int_{-\infty}^{x}f(t)dt \\
    &= \int_{0}^{x}\lambda e^{-\lambda t}dt \\
    &= -\int_{0}^{x}e^{u}du \\
    &=-\left(e^u\right)\Big|_{t=0}^x \\
    &=-\left(e^{-\lambda t}\right)\Big|_{t=0}^x \\
    &= 1-e^{\lambda x}.
\end{align*}
Luego, notemos que
\begin{align*}
    \P(X\leq P_n)=p &\iff F_X(P_n)=p \\
    &\iff 1-e^{-\lambda P_n}=p \\
    &\iff -\lambda P_n = \ln(1-p) \\
    &\iff P_n = -\frac{\ln(1-p)}{\lambda},
\end{align*}
obteniendo así el resultado buscado. \hfill $\blacktriangle$

\section*{Ejercicio 3.}
En ciertos sistemas, un cliente es asignado a una de dos prestadoras de servicio. Si el tiempo para que el cliente sea atendido por la prestadora de servicio $i$ tiene una distribución exponencial con parámetro $\lambda_i(i = 1, 2)$, y $p$ es la proporción de todos los clientes atendidos por la prestadora de servicio 1, entonces la función de densidad de probabilidad de
\[
X = "\text{tiempo para ser atendido de un cliente seleccionado al azar}"
\]
es
\[
f(x; \lambda_1, \lambda_2, p) =
\begin{cases}
    p\lambda_1e^{-\lambda_1x} + (1 - p) \lambda_2e^{-\lambda_2x}, & \text{si } x \geq 0,\\
    0, & \text{en otro caso}.
\end{cases}
\]
\begin{enumerate}
    \item[(i)] ¿Cuál es la función de distribución asociada a $X$?
    \item[(ii)] Determine $\E(X)$.
    \item[(iii)] Encuentre $\Var(X)$.
\end{enumerate}
\textbf{Solución.}
\begin{enumerate}
    \item[(i)] Tenemos que
        \begin{align*}
            F_X(x) &= \int_{-\infty}^{x}f(t)dt \\
            &= \int_{0}^{x}\left(p\lambda_1e^{-\lambda_1t} + (1 - p) \lambda_2e^{-\lambda_2t}\right)dt \\
            &= p\int_{0}^{x}\left(\lambda_1e^{-\lambda_1t}\right)dt + (1-p)\int_{0}^{x}\left(\lambda_2e^{-\lambda_2t}\right)dt \\
            &= p(1-e^{-\lambda_1 x}) + (1-p)(1-e^{-\lambda_2 x}) \\
            &= 1 - pe^{-\lambda_1 x} - (1-p)e^{-\lambda_2 x}.
        \end{align*}
    \item[(ii)] Sabemos que, para una variable aleatoria $Y$ exponencial de parámetro $\lambda > 0$, (con $u=-\lambda x$) se tiene que
        \begin{align*}
            \E(Y) &= \int_{-\infty}^{\infty}xf(x)dx \\
            &= \int_{0}^{\infty}x\lambda e^{-\lambda x}dx \\
            &= \frac{1}{\lambda}\int_{0}^{\infty}ue^udu \\
            &= \frac{1}{\lambda}\left(ue^u-e^u\right)\Big|_{x=0}^\infty \\
            &= \frac{1}{\lambda}\left(e^{-\lambda x}(-\lambda x-1)\right)\Big|_{x=0}^\infty \\
            &= \frac{1}{\lambda}.
        \end{align*}
        Así, afirmamos que, para la variable aleatoria $X$ original
        \begin{align*}
            \E(X) &= \int_{-\infty}^{\infty}xf(x)dx \\
            &= \int_{0}^{\infty}\left(p\lambda_1xe^{-\lambda_1x} + (1 - p) \lambda_2xe^{-\lambda_2x}\right)dx \\
            &= p\int_{0}^{\infty}\left(\lambda_1xe^{-\lambda_1x}\right)dx + (1-p)\int_{0}^{\infty}\left(\lambda_2xe^{-\lambda_2x}\right)dx \\
            &= \frac{p}{\lambda_1} + \frac{1-p}{\lambda_2}.
        \end{align*}
    \item[(iii)] Nuevamente, notemos que para $Y$ (usando nuevamente el cambio de variable $u=-\lambda x$)
        \begin{align*}
            \E[Y^2] &= \int_{-\infty}^{\infty}x^2f(x)dx \\
            &= \int_{0}^{\infty}x^2\lambda e^{-\lambda x}dx \\
            &= \frac{1}{\lambda}\int_{0}^{\infty}tue^udu \\
            &= -\frac{1}{\lambda^2}\int_{0}^{\infty}u^2e^udu \\
            &= -\frac{1}{\lambda^2}\left(u^2e^u+2(e^u(1-u))\right)\Big|_{x=0}^\infty \\
            &= -\frac{1}{\lambda^2}\left((-\lambda x)^2e^{-\lambda x}+2(e^{-\lambda x}(1+\lambda x))\right)\Big|_{x=0}^\infty \\
            &= \frac{2}{\lambda^2},
        \end{align*}
        y en consecuencia
        \begin{align*}
            \E(X^2) &= \int_{-\infty}^{\infty}x^2f(x)dx \\
            &= \int_{0}^{\infty}\left(p\lambda_1x^2e^{-\lambda_1x} + (1 - p) \lambda_2x^2e^{-\lambda_2x}\right)dx \\
            &= p\int_{0}^{\infty}\left(\lambda_1x^2e^{-\lambda_1x}\right)dx + (1-p)\int_{0}^{\infty}\left(\lambda_2x^2e^{-\lambda_2x}\right)dx \\
            &= \frac{2p}{\lambda_1^2} + \frac{2(1-p)}{\lambda_2^2}
        \end{align*}
        Luego, de lo anterior y de lo mencionado en el inciso $(ii)$ concluimos que
        \begin{align*}
            \text{Var}(X) &= \frac{2p}{\lambda_1^2} + \frac{2(1-p)}{\lambda_2^2} - \left(\frac{p}{\lambda_1} + \frac{1-p}{\lambda_2}\right)^2 \\
            &= \frac{2p}{\lambda_1^2} + \frac{2(1-p)}{\lambda_2^2} - \left(\frac{p^2}{\lambda_1^2}+\frac{2p(1-p)}{\lambda_1\lambda_2}+\frac{(1-p)^2}{\lambda_2^2}\right) \\
            &= \frac{2p-p^2}{\lambda_2} + \frac{2(1-p)-(1-p)^2}{\lambda_2} - \frac{2p(1-p)}{\lambda_1\lambda_2}.
        \end{align*}
        \hfill $\blacktriangle$
\end{enumerate}

\section*{Ejercicio 4.}
Sea $X$ una variable aleatoria con distribución de Weibull con parámetros $\alpha = 2$ y $\beta > 0$. Demuestre que
\[
Y = \dfrac{2X^2}{\beta^2}
\]
tiene una distribución ji-cuadrado con $\nu = 2$ grados de libertad.\\
\textbf{Demostración.} \\
Tenemos que $X$ tiene función de densidad
\[
f(t)=
\begin{cases}
    \frac{2}{\beta}\left(\frac{x}{\beta}\right)e^{-\left(x/\beta\right)^2}, &x \geq 0 \\
    0, &x<0
\end{cases}
\]
Así, tomando $u=-(x/\beta)^2$,
\begin{align*}
    F_X(x) &= \int_{-\infty}^{x}f(t)dt \\
    &= \int_{0}^{x}\frac{2}{\beta}\left(\frac{t}{\beta}\right)e^{-\left(t/\beta\right)^2}dt \\
    &= -\int_{0}^{x}e^udu \\
    &= -\left(e^u\right)\Big|_{t=0}^x \\
    &= -\left(e^{-(x/\beta)^2}\right)\Big|_{t=0}^x \\
    &= 1-e^{-(x/\beta)^2}.
\end{align*}
Notemos que
\begin{align*}
    F_Y(y) &= \P(Y\leq y) \\
    &= \P\left(\frac{2X^2}{\beta^2}\leq y\right) \\
    &= \P\left(X\leq \frac{\beta\sqrt{y}}{\sqrt{2}}\right) \\
    &= F_X\left(\frac{\beta\sqrt{y}}{\sqrt{2}}\right) \\
    &= 1-e^{\left(-\beta\sqrt{y}/(\beta\sqrt{2})\right)^2} \\
    &= 1-e^{-y/2},
\end{align*}
que coincide justamente con la función de distribución de una variable aleatoria con distribución ji-cuadrado de parámetros $\nu=2$,
de lo que se sigue el resultado. \hfill $\blacktriangle$

\section*{Ejercicio 5.}
Sean $X$ y $Y$ dos variables aleatorias con valores en el conjunto $\llaves{x_1, \dots, x_n}$, con funciones de probabilidad (masas de probabilidad) $f_X$ y $f_Y$, respectivamente. Demuestre que
\[
- \E ( \log f_Y (X)) \geq - \E ( \log f_X (X)),
\]
y que
\[
- \E( \log f_X (X)) \geq \log n.
\]
¿Cuándo se tienen las igualdades?\\
La función
\[
H(X) := - \E (\log f_X (X))
\]
se llama \textit{entropía} de la variable aleatoria $X$. La primera desigualdad se cumple de manera mucho mas general: si $Y := g(X)$ para alguna función $g$, entonces $H(Y) \leq H(X)$.\\
\textbf{\textit{Demostración.}}

\end{document}